\documentclass[12pt,a4paper]{article}
\usepackage{fontspec}
\setmainfont{DejaVu Serif}
\setsansfont{DejaVu Sans}
\usepackage[brazil]{babel}
\usepackage[margin=2.2cm]{geometry}
\usepackage{booktabs,tabularx,array,longtable}
\usepackage{xcolor}
\usepackage{parskip}
\usepackage{amsmath}
\usepackage{enumitem}
\usepackage{fancyhdr}
\setlength{\headheight}{15pt}
\pagestyle{fancy}
\fancyhf{}
\lhead{Análise da Complexidade do Sistema Gérard}
\rhead{Atualização C67}
\cfoot{\thepage}
\definecolor{vermelhosuave}{RGB}{194,104,104}
\begin{document}

\section*{Atualização da análise de complexidade - ciclos C36 a C67}

Esta atualização amplia a análise anterior da comparação de medidas e incorpora a generalização arquitetural do ciclo C67. O sistema passou a tratar os diagramas de Vergnaud, os diagramas complementares, as barras, os cartões e os eixos como projeções manipuláveis de um \textbf{estado semântico compartilhado}. Uma ação deixa de atualizar diretamente outra representação: ela altera o estado do domínio, que resolve a relação matemática e redistribui um mesmo snapshot para todas as visualizações.

\section{Síntese do impacto}

A complexidade tecnológica geral permanece baixa a média, mas a complexidade funcional e de domínio permanece alta. O ciclo C67 reduz a complexidade acidental provocada por cadeias locais de sincronização, ao custo de introduzir uma camada explícita de coordenação.

\begin{tabularx}{\textwidth}{p{3.8cm}p{2.8cm}X}
\toprule
\textbf{Dimensão} & \textbf{Avaliação} & \textbf{Impacto acumulado até C67}\\
\midrule
Complexidade de domínio & Alta & Papéis semânticos, sinais, incógnita e relações devem permanecer coerentes entre representações distintas.\\
Complexidade de interação & Alta & Arraste, edição, exclusão, eixos e protocolos podem iniciar a mesma transação semântica.\\
Complexidade arquitetural & Média para alta & Foi criado um estado compartilhado com snapshot, origem da ação, papel alterado e proteção contra recursão.\\
Complexidade de renderização & Baixa a média & Coleções e barras são reconstruídas a partir das quantidades consolidadas.\\
Complexidade de manutenção & Média & Invariantes passam a estar centralizados, reduzindo correções duplicadas em componentes visuais.\\
\bottomrule
\end{tabularx}

\section{Modelo semântico compartilhado}

Para os três papéis principais das categorias atuais, o estado é representado por:
\[
S = (T, V, K, i, o, \nu),
\]
em que $T$ é a categoria, $V=(v_1,v_2,v_3)$ contém os valores, $K=(k_1,k_2,k_3)$ indica quais valores foram efetivamente modelados, $i$ identifica o papel alterado, $o$ registra a origem da ação e $\nu$ é a versão do snapshot.

A relação canônica utilizada no nível dos três papéis é:
\[
v_1 + v_2 = v_3.
\]

Essa forma cobre, entre outros casos:
\begin{itemize}[noitemsep]
  \item Parte 1 $+$ Parte 2 $=$ Todo;
  \item Estado inicial $+$ Transformação $=$ Estado final;
  \item Referido $+$ Valor relativo $=$ Referendo;
  \item Relação inicial $+$ Transformação $=$ Relação final;
  \item Transformação 1 $+$ Transformação 2 $=$ Transformação resultante.
\end{itemize}

O papel alterado define a direção da resolução. Se $v_3$ é alterado e $v_1$ é conhecido, então $v_2=v_3-v_1$. Se $v_1$ ou $v_2$ é alterado e ambos são conhecidos, então $v_3=v_1+v_2$. Quando não há papel explicitamente alterado, somente um termo ausente é derivado.

\section{Fluxo reativo}

O fluxo consolidado é:

\begin{center}
Ação do usuário $\rightarrow$ captura dos valores $\rightarrow$ estado compartilhado $\rightarrow$ resolução $A+B=C$ $\rightarrow$ snapshot $\rightarrow$ atualização de todas as representações.
\end{center}

As origens previstas são Vergnaud, diagrama complementar, eixo x, eixo vertical, edição textual, arraste, exclusão e protocolo de scaffolding. Uma trava booleana de aplicação impede que o redesenho provocado pelo snapshot seja interpretado como uma nova ação. Sem essa trava, uma atualização Vergnaud $\rightarrow$ Venn poderia acionar novamente Venn $\rightarrow$ Vergnaud, produzindo recursão ou oscilações.

\section{Complexidade temporal}

Considere $p=3$ como o número de papéis principais, $r$ como o número de representações coordenadas e $q$ como o total de unidades gráficas manipuláveis.

\subsection{Captura e resolução}

A leitura dos três elementos de Vergnaud e a resolução da relação são constantes:
\[
O(p)=O(1).
\]

No diagrama complementar, a contagem dos quadradinhos percorre as unidades. Como há número constante de regiões, o custo é:
\[
O(q).
\]

A criação do snapshot e a escolha do termo dependente permanecem $O(1)$.

\subsection{Aplicação nas representações}

Atualizar os três valores textuais de Vergnaud custa $O(1)$. Reconstruir as coleções ou barras custa $O(q')$, em que $q'$ é o número de unidades do novo estado. Atualizar os eixos e cartões é $O(1)$.

Na comparação de medidas, a correspondência cromática ainda exige ordenar as unidades de cada barra de baixo para cima. Para $q=r_f+r_e$:
\[
O(r_f\log r_f+r_e\log r_e)=O(q\log q).
\]

Logo, o custo total de uma transação é:
\[
O(q) \quad \text{nas coleções simples},
\]
e
\[
O(q\log q) \quad \text{na comparação com pareamento cromático}.
\]

Como as quantidades educacionais são pequenas e limitadas pela área de desenho, esse custo não representa gargalo prático.

\section{Complexidade espacial}

O estado compartilhado armazena apenas três valores, três indicadores e metadados constantes, portanto utiliza:
\[
O(1).
\]

As unidades gráficas continuam dominando a memória:
\[
O(q).
\]

A centralização não duplica os quadradinhos no modelo de domínio. O snapshot contém quantidades, não cópias das figuras.

\section{Invariantes de consistência}

Após qualquer transação, devem ser verdadeiras as seguintes condições:
\begin{enumerate}[noitemsep]
  \item todas as representações exibem o mesmo valor para cada papel conhecido;
  \item a relação $v_1+v_2=v_3$ é preservada quando há informação suficiente;
  \item o sinal do segundo termo é preservado nos números relativos;
  \item a incógnita permanece desconhecida até ser modelada ou derivada pela relação;
  \item eixos, cartões e pontos de controle refletem o mesmo valor da relação;
  \item quadradinhos e rótulos numéricos representam a mesma quantidade;
  \item uma atualização de renderização não cria uma segunda ação do usuário.
\end{enumerate}

\section{Impacto sobre a manutenção}

Antes do C67, correções eram adicionadas em pares: eixo vertical para barras, barras para Vergnaud, cartão para eixo, Venn para valores e valores para Venn. Esse desenho aumenta o número potencial de conexões entre $r$ representações para até:
\[
O(r^2).
\]

Com o estado compartilhado, cada representação se conecta ao coordenador, reduzindo a topologia conceitual para:
\[
O(r).
\]

O ganho é de manutenção, não necessariamente de tempo de execução. Novas representações precisam implementar duas operações: publicar alterações no estado e renderizar um snapshot. Elas não precisam conhecer todas as demais visualizações.

\section{Riscos e controles}

Os principais riscos introduzidos pela centralização são:
\begin{itemize}[noitemsep]
  \item derivar um termo que o desenho pedagógico deveria manter vazio;
  \item escolher o termo dependente errado quando duas regiões são alteradas na mesma ação;
  \item reconstruir componentes durante um arraste contínuo em frequência excessiva;
  \item apagar a distinção entre valor curado e valor efetivamente modelado pelo usuário.
\end{itemize}

Os controles implementados incluem indicadores de valor conhecido, registro do papel alterado, origem da ação, inicialização vazia e trava contra recursão. A validação deve continuar incluindo inspeção visual por categoria, além dos testes do modelo.

\section{Verificação do ciclo C67}

Foram executadas:
\begin{itemize}[noitemsep]
  \item compilação completa pelo Ant/NetBeans;
  \item regressão estrutural, de internacionalização, vínculos, curadoria e logs;
  \item teste isolado do estado compartilhado;
  \item casos de composição $8+6=14$, alteração do Todo, comparação $6+8=14$ e transformação $10+(-3)=7$.
\end{itemize}

A inspeção visual exaustiva de todas as combinações de arraste permanece necessária, pois testes unitários do estado não verificam sobreposição, posicionamento ou fluidez da interface.

\section{Classificação do ciclo C67}

O ciclo C67 é classificado como \textbf{estrutural e de generalização}. Ele não corrige apenas uma figura específica. Formaliza uma regra arquitetural válida para todas as categorias: representações manipuláveis são projeções de um mesmo estado matemático e devem ser atualizadas consistentemente após toda ação.

\section{Conclusão}

A introdução do estado semântico compartilhado aumenta a formalização do domínio, mas reduz a complexidade acidental das sincronizações locais. O custo assintótico permanece dominado pela quantidade de unidades gráficas, enquanto o estado e a resolução da relação são constantes. O principal benefício é a coerência: Vergnaud, Venn, barras, cartões e eixos deixam de competir como fontes de verdade e passam a representar o mesmo snapshot matemático.

\end{document}
